\chapter{Связывающий родитель Мориты и носитель переходов размерности 300}

Угловое чтение уже переносит скалярную семейную кривизну в выбранное
состояние. Двусторонняя связность нужна теперь только для динамики
взаимодействий. Прежде чем задавать её компоненты, определим сам бимодуль
между двумя чтениями.

\section{Антикруговая сверка}

Эта глава не повторяет отображающий конус, фактор форм Конна или квадрат
стандартной суперсвязности. Там вычислялись разные абсолютные кривизны на
одном конечном графе. Здесь исследуется пространство переходов между двумя
диагональными чтениями общего носителя.

Связывающие алгебры и эквивалентность Мориты являются стандартными
объектами операторной теории. Двусторонние связности и их кривизны также
известны; см. \path{arXiv:q-alg/9512004}, \path{arXiv:0904.0539},
\path{arXiv:1911.05008} и \path{arXiv:2302.00489}. Проектный результат
может состоять только в точном совпадении этих структур с независимо
полученными размерностями, чтениями и следами.

\section{Триста состояний как прямоугольный бимодуль}

Размерности семейного и наблюдаемого чтений равны двадцати и пятнадцати.
Поэтому
\begin{equation}
 H_{300}\simeq E:=M_{20\times15}(\mathbb C).
 \label{eq:v5-morita-carrier}
\end{equation}
Алгебра $M_{20}(\mathbb C)$ действует на $E$ слева, а
$M_{15}(\mathbb C)$ --- справа. Эти действия коммутируют и воспроизводят
два факторных чтения прежнего коммутирующего квадрата.

Теперь различаются два естественных объекта:
\begin{equation}
 \operatorname{End}_{\mathbb C}(E)\simeq M_{300}(\mathbb C),
 \qquad
 \mathcal L(E)\simeq M_{35}(\mathbb C).
 \label{eq:v5-endomorphism-linking-distinction}
\end{equation}
Первый является алгеброй операторов на переходах. Второй собирает два
чтения и сами переходы в одной блочной алгебре.

\section{Связывающая алгебра и единственный след}

Имеем
\begin{equation}
 \mathcal L(E)=
 \begin{pmatrix}
  M_{20}(\mathbb C)&E\\
  E^*&M_{15}(\mathbb C)
 \end{pmatrix}
 \simeq M_{35}(\mathbb C).
 \label{eq:v5-morita-linking-algebra}
\end{equation}
Размерностная ведомость замыкается точно:
\begin{equation}
 20^2+15^2+2\cdot20\cdot15
 =400+225+300+300=35^2.
\end{equation}
Следовательно, исходный носитель размерности 300 является ровно одним
внедиагональным углом связывающей алгебры.

Единственный нормированный след $\tau_{35}$ даёт
\begin{equation}
 \tau_{35}(p_{20})=\frac47,
 \qquad
 \tau_{35}(p_{15})=\frac37.
\end{equation}
После условной нормировки углов получаются прежние следы:
\begin{equation}
 \frac{\tau_{35}(p_{20}Ap_{20})}{\tau_{35}(p_{20})}=\tau_{20}(A),
 \qquad
 \frac{\tau_{35}(p_{15}Bp_{15})}{\tau_{35}(p_{15})}=\tau_{15}(B).
 \label{eq:v5-conditioned-linking-traces}
\end{equation}
Числа $4/7$ и $3/7$ следуют из рангов углов, а не являются секторными
параметрами.

Точное наблюдаемое разложение $15=6+2+3+3+1$ индуцирует
\begin{equation}
 300=120+40+60+60+20.
 \label{eq:v5-observed-linking-blocks}
\end{equation}
Так пять наблюдаемых блоков входят в переходный носитель без добавления
новых состояний.

\section{Двусторонняя связность}

Пусть $A_F$ и $A_O$ --- связности двух диагональных чтений. На
$\xi\in E$ ковариантная производная имеет вид
\begin{equation}
 \nabla_E\xi=d\xi+A_F\xi-\xi A_O.
 \label{eq:v5-morita-bimodule-connection}
\end{equation}
Знак минус следует из правого действия и закона преобразования
\begin{equation}
 \xi\longmapsto U_F\xi U_O^*.
\end{equation}
Правое правило Лейбница фиксируется обычным блочным умножением внутри
$M_{35}$; отдельное переставление форм не выбирается.

Кривизна переходов равна
\begin{equation}
 \mathcal R_E(\xi)=F_F\xi-\xi F_O.
 \label{eq:v5-morita-relative-curvature}
\end{equation}
На векторизации $E$ левое и правое действия представлены операторами
\begin{equation}
 L_{F_F}=F_F\otimes I_{15},
 \qquad
 R_{F_O}=I_{20}\otimes F_O^T.
\end{equation}
Поэтому
\begin{align}
 \tau_{300}\!\left((L_{F_F}-R_{F_O})^2\right)
 &=\tau_{20}(F_F^2)+\tau_{15}(F_O^2)\nonumber\\
 &\quad-2\tau_{20}(F_F)\tau_{15}(F_O).
 \label{eq:v5-morita-relative-trace}
\end{align}
Для центрированных кривизн смешанный член исчезает. Ранее найденное
ортогональное разложение одного следа тем самым является нормой
относительной кривизны бимодуля.

Машинный дефект общей следовой формулы меньше $7\cdot10^{-19}$,
центрированной формулы --- меньше $8\cdot10^{-15}$, а ковариантности ---
меньше $1.2\cdot10^{-13}$.

\begin{theorem}[Связывающий родитель двух чтений]
Носитель $H_{300}$ канонически является бимодулем
$M_{20\times15}(\mathbb C)$. Его связывающая алгебра $M_{35}(\mathbb C)$
содержит два чтения в диагональных углах, а все 300 переходных направлений
--- во внедиагональном угле. Единственный след восстанавливает оба
нормированных угловых следа, а двусторонняя кривизна имеет канонический знак
разности.
\end{theorem}

\section{Калибровочная граница}

$M_{35}(\mathbb C)$ нельзя объявлять новой координатной алгеброй. Полная
унитарная алгебра Ли имела бы размерность $35^2=1225$, то есть на 1213
направлений больше калибровочной алгебры Стандартной модели. Связывающая
алгебра является операторным контейнером, а не автоматически физической
алгеброй координат.

Правый угол $M_{15}$ пока остаётся полным контрольным чтением. Его нужно
ограничить точным действием
$\mathbb C\oplus\mathbb H\oplus M_3(\mathbb C)$, а левый угол --- уже
полученным аффинным семейным действием. Только после этого допустимо
классифицировать компоненты $\nabla_E$ и проверять происхождение юкавского
оператора.

\begin{nogocriterion}[Запрет калибровки связывающего контейнера]
Если все внутренние унитарные преобразования $M_{35}$ объявляются
физической калибровочной группой, путь закрывается. Допустима только
связность, совместимая с уже зафиксированными семейным и
стандарт-модельным действиями.
\end{nogocriterion}

\begin{protocol}
\begin{enumerate}
 \item происхождение размерности 300 как пространства переходов:
 \textbf{получено};
 \item связывающая алгебра $M_{35}$: \textbf{получена};
 \item различение ролей $M_{35}$ и $M_{300}$: \textbf{получено};
 \item один след и нормированные угловые следы: \textbf{получены};
 \item двусторонняя кривизна без свободного переставления:
 \textbf{получена};
 \item точное наблюдаемое разложение переходов: \textbf{получено};
 \item $M_{35}$ как координатная алгебра: \textbf{запрещено};
 \item юкавский оператор как единственная компонента связности:
 \textbf{открыто};
 \item полный БВ/БРСТ-гессиан: \textbf{открыт};
 \item физическое замыкание: \textbf{не достигнуто}.
\end{enumerate}
\end{protocol}

Следующий гейт должен ограничить диагональные углы известными физическими
действиями и классифицировать совместимые связности на $E$. Главный вопрос:
становится ли ранг-один юкавский угол единственной внедиагональной
компонентой или для его выбора вновь требуется внешний селектор.

Машинный сертификат сохранён в
\path{s2t_v5_morita_linking_parent_gate.py} и
\path{s2t_v5_morita_linking_parent_gate_results.json}.